\section{当前实现的问题分析}

\begin{enumerate}[itemsep=0.55em]
  \item \textbf{83 维输入中有 3 维可重建。}
  由式~\eqref{eq:redundant}，常规控制帧中，最新观测的前 3 维等于 \code{cmd_ori} 中两组三维向量的和。保留原始指令和上一帧修正，即可从 MLP 输入中省去合成指令三维。去冗余对训练性能的影响仍需消融验证。\ev{C2,C4}

  \item \textbf{仿真施力与头部附近牵引存在差异。}
  仿真在 \code{base} 刚体质心施力，附加力矩为零。实物若从头部附近牵引，拉点相对质心的偏移会产生等效力矩，见式~\eqref{eq:force-transform}。基线每隔 $0.02\,\mathrm{s}$ 只施力 $0.005\,\mathrm{s}$，参考映射仍使用脉冲的完整幅值。\ev{C7}

  \item \textbf{平移参考修正越过阈值就饱和。}
  式~\eqref{eq:force-map} 中，任一平移轴的力幅值刚超过 $17.5\,\mathrm{N}$，限幅前的修正幅值就已超过 $1.75\,\mathrm{m/s}$，直接达到该轴的限幅。例如，$18\,\mathrm{N}$ 和 $50\,\mathrm{N}$ 的前向力都对应 $1.2\,\mathrm{m/s}$ 的参考修正，阈值以上的力幅值差异无法反映到平移参考中。\ev{C3,C5}

  \item \textbf{运动跟踪奖励的设计与牵引任务不匹配。}
  底层负责跟踪给定速度，牵引 MLP 从观测中学习与外力对应的速度修正。当前平移和偏航跟踪奖励仍计入牵引 PPO 回报，而底层已冻结。其跟踪目标为原始指令加上一帧 MLP 修正，策略可能通过选取底层更容易执行的修正来提高跟踪分，即使该修正偏离外力参考。这两项跟踪奖励应留在底层控制训练中。\ev{C6,C8}

  \item \textbf{撤力后的实际减速过程与停车判据未明确。}
  当前奖励未明确约束实际速度的减速时间和平滑程度。原始指令为零时，应以实际平移速度与偏航角速度在限时内降至各自阈值以下、并持续一段时间作为停车判据；原始指令非零时，应评价恢复跟踪该指令的过程。\ev{C5,C6,C7}

  \item \textbf{力标签比实际施力提前一帧。}
  当前控制步先用旧力完成仿真，再更新力标签，并用新标签计算本帧的 \code{compliance}。例如，撤力边界这一帧仍受最后一次非零力脉冲，奖励目标却已归零，到下一帧才停止施力。这是切换时相差一个控制步（$0.02\,\mathrm{s}$）的时序错误，基线因复现原实现而保留了这一行为。\ev{C6,C7,C8}
\end{enumerate}
